\chapter*{Цели и задачи}
\addcontentsline{toc}{chapter}{Введение}
\label{ch:intro_lab4}

\textbf{Цель лабораторной работы:} изучение принципов построения информационных систем с использованием логических методов классификации.

\textbf{Основные задачи:}
\begin{itemize}
    \item освоение технологии внедрения алгоритмов на основе решающих списков в приложения;
    \item освоение технологии внедрения алгоритмов на основе решающих деревьев в приложения;
    \item изучение параметров логической классификации;
    \item освоение модификаций логических методов классификации.
\end{itemize}

В качестве набора данных для выполнения индивидуального задания выбран набор данных \textbf{Abalone} (морское ушко), который содержит информацию о возрасте моллюсков, определяемом по количеству колец на раковине. Исходная задача регрессии преобразована в задачу классификации на три возрастные категории: молодые (0-9 колец), взрослые (10-15 колец) и старые (16 и более колец).
\endinput
